%! Unidad 7 — La química del carbono: ejercicios resueltos, ejercicios de aplicación, de ampliación y autoevaluación
%! Fuente: Fuentes/Quimica/Unidad 7 («Ejercicios resueltos», «Practica y calcula» y «Autoevaluación», capturas 570-572).
%! Enunciados adaptados con compuestos y datos cambiados; soluciones rehechas. La fuente no trae olimpiadas en este tema.
%! Dificultad: \(\star\) baja, \(\star\star\) media, \(\star\star\star\) alta.

\sloppy % nombres largos sin guiones (\fussy al final)

\section{Ejercicios resueltos}\label{sec:ud7-resueltos}

\begin{example}
    \textbf{Nombrar un compuesto con más de un grupo funcional.} Nombra \chemfig{CH_3-CH(-[6]CH_3)-CH(-[6]OH)-CH=CH-COOH}.

    \textbf{Solución.} \textbf{1. Grupos funcionales}: \ch{-COOH} y \ch{-OH}. El ácido tiene más prioridad: se nombra con el sufijo
    \emph{-oico} y el alcohol, con el prefijo \emph{hidroxi-}.

    \textbf{2. Cadena principal}: la que contiene el \ch{-COOH} y es más larga: 6 carbonos, con el C1 en el \ch{-COOH}. Así, el
    doble enlace queda entre C2 y C3, el \ch{-OH} en C4 y el metilo en C5.

    \textbf{3. Sustituyentes por orden alfabético}: 4-hidroxi y 5-metil. \textbf{4. Cadena}: hex-2-en- + \emph{-oico}.

    El nombre es \textbf{ácido 4-hidroxi-5-metilhex-2-enoico}.
\end{example}

\begin{example}
    \textbf{Formular un compuesto con más de un grupo funcional.} Formula el 4-amino-3-oxopentanal.

    \textbf{Solución.} El sufijo \emph{-al} indica que el grupo principal es el aldehído, en el C1 de una cadena de 5 carbonos
    (\emph{pent-}). Los prefijos indican una cetona (\emph{oxo-}) en el C3 y un grupo amino en el C4. Dibujamos la cadena, colocamos
    los grupos y completamos con hidrógenos:
    \[ \chemfig{CH_3-CH(-[6]NH_2)-CO-CH_2-CHO} \]
\end{example}

\begin{example}
    \textbf{Fórmula de una amina e isómeros.} Al quemar \qty{0.59}{\gram} de un compuesto de C, H y N se obtienen
    \qty{1.32}{\gram} de \ch{CO2} y \qty{0.81}{\gram} de \ch{H2O}. Su masa molar es \qty{59}{\gram\per\mol}. a) Halla su fórmula
    molecular. b) Escribe y nombra los isómeros posibles si es una amina.

    \textbf{Solución.} a) \(m(\ch{C}) = 1{,}32 \cdot \frac{12}{44} = \qty{0.36}{\gram}\) (\qty{0.030}{\mol});
    \(m(\ch{H}) = 0{,}81 \cdot \frac{2}{18} = \qty{0.090}{\gram}\) (\qty{0.090}{\mol}); \(m(\ch{N}) = 0{,}59 - 0{,}36 - 0{,}090 =
    \qty{0.14}{\gram}\) (\qty{0.010}{\mol}). Proporción \(3 : 9 : 1\): empírica \ch{C3H9N}, de masa 59, igual que la molar. La
    fórmula molecular es \ch{C3H9N}.

    b) Primarias: propan-1-amina, \ch{CH3-CH2-CH2-NH2}, y propan-2-amina, \ch{CH3-CH(NH2)-CH3}. Secundaria: N-metiletanamina,
    \ch{CH3-NH-CH2-CH3}. Terciaria: N,N-dimetilmetanamina (trimetilamina), \ch{N(CH3)3}.
\end{example}

\begin{example}
    \textbf{Todos los isómeros de una fórmula.} Escribe y nombra todos los hidrocarburos de fórmula \ch{C4H8}, incluidos los
    cíclicos y los estereoisómeros.

    \textbf{Solución.} La fórmula \(\ch{C_nH_{2n}}\) corresponde a un alqueno o a un cicloalcano.
    \begin{itemize}
        \item Alquenos de cadena lineal: but-1-eno, \ch{CH2=CH-CH2-CH3}, y but-2-eno, \ch{CH3-CH=CH-CH3}, que tiene dos isómeros
        geométricos, \emph{cis} y \emph{trans}.
        \item Alqueno ramificado: metilpropeno, \chemfig{CH_2=C(-[6]CH_3)-CH_3}.
        \item Cíclicos: ciclobutano y metilciclopropano:
        \[ \chemfig{*4(----)} \qquad\qquad \chemfig{*3(--(-CH_3)-)} \]
    \end{itemize}
    En total, seis compuestos distintos.
\end{example}

\begin{example}
    \textbf{Estequiometría de una esterificación.} Se calientan \qty{30}{\gram} de ácido acético con exceso de etanol y un poco de
    ácido sulfúrico como catalizador. Si el rendimiento es del 70\,\%, ¿qué masa de etanoato de etilo se obtiene?

    \textbf{Solución.}
    \[ \ch{CH3-COOH + CH3-CH2OH <=> CH3-COO-CH2-CH3 + H2O} \]
    \(M(\ch{CH3COOH}) = \qty{60}{\gram\per\mol}\) y
    \(M(\ch{CH3COOC2H5}) = \qty{88}{\gram\per\mol}\).
    \[ n(\text{ácido}) = \frac{30}{60} = \qty{0.50}{\mol} \;\Rightarrow\; n(\text{éster}) = 0{,}50 \cdot 0{,}70 = \qty{0.35}{\mol}
       \;\Rightarrow\; m = 0{,}35 \cdot 88 = \qty{30.8}{\gram} \]
\end{example}

%========================================================================================
\section{Ejercicios de aplicación}\label{sec:ud7-aplicacion}
%========================================================================================

\subsection{Fórmulas de los compuestos del carbono}

\begin{exercise}
    \(\star\star\) Al quemar completamente un hidrocarburo se obtienen \qty{33.0}{\gram} de \ch{CO2} y \qty{13.5}{\gram} de
    \ch{H2O}. Determina su fórmula empírica y, si su masa molar es \qty{84}{\gram\per\mol}, su fórmula molecular.
    \begin{solu}
        \textcolor{blue}{\(n(\ch{C}) = 33{,}0/44 = \qty{0.75}{\mol}\); \(n(\ch{H}) = 2 \cdot 13{,}5/18 = \qty{1.5}{\mol}\); proporción
        1 : 2.} \textcolor{red}{\ch{CH2}}; \(84/14 = 6\): \textcolor{red}{\ch{C6H12}}.
    \end{solu}
\end{exercise}

\begin{exercise}
    \(\star\star\star\) Al quemar \qty{2.15}{\gram} de un compuesto de C, H y O se obtienen \qty{4.40}{\gram} de \ch{CO2} y
    \qty{1.35}{\gram} de \ch{H2O}. Si \qty{0.43}{\gram} de su vapor ocupan \qty{161}{\mL} a \qty{120}{\degreeCelsius} y
    \qty{1}{\atm}, halla sus fórmulas empírica y molecular.
    \begin{solu}
        \textcolor{blue}{\(m(\ch{C}) = 4{,}40 \cdot 12/44 = \qty{1.20}{\gram}\) (\qty{0.100}{\mol}); \(m(\ch{H}) = 1{,}35 \cdot 2/18 =
        \qty{0.15}{\gram}\) (\qty{0.150}{\mol}); \(m(\ch{O}) = 2{,}15 - 1{,}35 = \qty{0.80}{\gram}\) (\qty{0.050}{\mol}). Proporción
        2 : 3 : 1.} Empírica: \textcolor{red}{\ch{C2H3O}} (43). \textcolor{blue}{\(M = \dfrac{0{,}43 \cdot 0{,}082 \cdot 393}{1 \cdot
        0{,}161} = \qty{86}{\gram\per\mol}\); \(86/43 = 2\).} Molecular: \textcolor{red}{\ch{C4H6O2}}.
    \end{solu}
\end{exercise}

\begin{exercise}
    \(\star\star\) Un compuesto contiene un 40,0\,\% de C, un 6,7\,\% de H y el resto de O. Su masa molar es
    \qty{60}{\gram\per\mol}. Halla su fórmula molecular y escribe dos compuestos que puedan tenerla.
    \begin{solu}
        \textcolor{blue}{En \qty{100}{\gram}: \(40{,}0/12 = 3{,}33\) mol de C; \(6{,}7\) mol de H; \(53{,}3/16 = 3{,}33\) mol de O:
        \ch{CH2O} (30); \(60/30 = 2\).} \textcolor{red}{\ch{C2H4O2}}: \textcolor{red}{ácido etanoico, \ch{CH3-COOH}, y metanoato de
        metilo, \ch{HCOO-CH3}}.
    \end{solu}
\end{exercise}

\begin{exercise}
    \(\star\star\) Escribe la fórmula molecular, la desarrollada, la semidesarrollada y la de esqueleto de: a) butano; b) propeno;
    c) propino; d) ciclobutano.
    \begin{solu}
        a) \textcolor{red}{\ch{C4H10}; \chemfig{H-C(-[2]H)(-[6]H)-C(-[2]H)(-[6]H)-C(-[2]H)(-[6]H)-C(-[2]H)(-[6]H)-H};
        \ch{CH3-CH2-CH2-CH3}; \chemfig{-[:30]-[:-30]-[:30]}}.
        b) \textcolor{red}{\ch{C3H6}; \chemfig{C(-[3]H)(-[5]H)=C(-[6]H)-C(-[2]H)(-[6]H)-H}; \ch{CH2=CH-CH3}; \chemfig{=[:30]-[:-30]}}.
        c) \textcolor{red}{\ch{C3H4}; \chemfig{H-C~C-C(-[2]H)(-[6]H)-H}; \ch{CH+C-CH3}; \chemfig{~-}}.
        d) \textcolor{red}{\ch{C4H8}; semidesarrollada \chemfig{CH_2*4(-CH_2-CH_2-CH_2-)} (en la desarrollada, cada \ch{CH2} se
        dibuja con sus dos enlaces \ch{C-H}); esqueleto \chemfig{*4(----)}}.
    \end{solu}
\end{exercise}

\begin{exercise}
    \(\star\) ¿Qué compuestos tienen la misma fórmula empírica? \ch{C2H2}, \ch{C2H4}, \ch{C6H6}, \ch{C3H6}, \ch{C4H10},
    \ch{C8H18}.
    \begin{solu}
        \textcolor{red}{\ch{C2H2} y \ch{C6H6} (\ch{CH}); \ch{C2H4} y \ch{C3H6} (\ch{CH2})}. El \ch{C4H10} es \ch{C2H5} y el
        \ch{C8H18}, \ch{C4H9}.
    \end{solu}
\end{exercise}

\subsection{Generalidades}

\begin{exercise}
    \(\star\star\) Localiza la cadena principal y los sustituyentes y nombra:
    \begin{gather*}
        \text{a)}\ \chemfig{CH_3-CH_2-CH(-[2]CH_2-[2]CH_2-[2]CH_3)-CH(-[6]CH_3)-CH_2-CH_3} \qquad
        \text{b)}\ \chemfig{CH_3-CH(-[6]CH_3)-CH_2-CH(-[2]CH_2-[0]CH_3)-CH(-[6]CH_3)-CH_3} \\
        \text{c)}\ \chemfig{CH_3-C(-[2]CH_3)(-[6]CH_3)-CH_2-C(-[2]CH_3)(-[6]CH_2-[0]CH_3)-CH_2-CH_3}
    \end{gather*}
    \begin{solu}
        a) \textcolor{blue}{La cadena más larga (7 C) entra por el propilo; el etilo de la izquierda es un sustituyente.}
        \textcolor{red}{4-etil-3-metilheptano}. b) \textcolor{blue}{Cadena horizontal de 6 C con tres ramificaciones; desde la
        derecha, \{2, 3, 5\}.} \textcolor{red}{3-etil-2,5-dimetilhexano}. c) \textcolor{blue}{Cadena de 6 C con cuatro
        ramificaciones, \{2, 2, 4, 4\}.} \textcolor{red}{4-etil-2,2,4-trimetilhexano}.
    \end{solu}
\end{exercise}

\begin{exercise}
    \(\star\star\) Escribe la fórmula semidesarrollada y el nombre de un compuesto de 5 carbonos que contenga: a) un grupo
    carbonilo en un extremo de la cadena; b) un grupo carboxilo; c) un grupo carbonilo en el interior de la cadena; d) un grupo
    amida; e) un grupo éster.
    \begin{solu}
        Por ejemplo: a) \textcolor{red}{\ch{CH3-(CH2)3-CHO}, pentanal}; b) \textcolor{red}{\ch{CH3-(CH2)3-COOH}, ácido
        pentanoico}; c) \textcolor{red}{\ch{CH3-CO-CH2-CH2-CH3}, pentan-2-ona}; d) \textcolor{red}{\ch{CH3-(CH2)3-CONH2},
        pentanamida}; e) \textcolor{red}{\ch{CH3-COO-CH2-CH2-CH3}, etanoato de propilo}.
    \end{solu}
\end{exercise}

\begin{exercise}
    \(\star\star\) Identifica los grupos funcionales y nombra: a) \ch{CH3-CHOH-CH2-CH2-CHO}; b) \ch{HOOC-CH2-CH3};
    c) \ch{CH3-CH2-CH2-NH-CH3}; d) \ch{CH3-CO-CH2-COO-CH3}.
    \begin{solu}
        a) \textcolor{red}{Aldehído (principal) y alcohol: 4-hidroxipentanal}. b) \textcolor{red}{Ácido carboxílico: ácido
        propanoico}. c) \textcolor{red}{Amina secundaria: N-metilpropan-1-amina}. d) \textcolor{red}{Éster (principal) y cetona:
        3-oxobutanoato de metilo}.
    \end{solu}
\end{exercise}

\subsection{Hidrocarburos}

\begin{exercise}
    \(\star\) Formula o nombra: a) heptano; b) 2,3,4-trimetilhexano; c) ciclopentano; d) 3-etil-2-metilpentano;
    \[ \text{e)}\ \chemfig{CH_3-CH(-[6]CH_3)-CH(-[6]CH_3)-C(-[2]CH_3)(-[6]CH_3)-CH_3} \qquad\qquad
       \text{f)}\ \chemfig{*4(---(-CH_2-CH_3)-)} \]
    \begin{solu}
        a) \textcolor{red}{\ch{CH3-(CH2)5-CH3}}; b) \textcolor{red}{\chemfig{CH_3-CH(-[6]CH_3)-CH(-[6]CH_3)-CH(-[6]CH_3)-CH_2-CH_3}};
        c) \textcolor{red}{\chemfig{*5(-----)}}; d) \textcolor{red}{\chemfig{CH_3-CH(-[6]CH_3)-CH(-[6]CH_2-[0]CH_3)-CH_2-CH_3}};
        e) \textcolor{blue}{\{2, 2, 3, 4\} desde la derecha frente a \{2, 3, 4, 4\}:} \textcolor{red}{2,2,3,4-tetrametilpentano};
        f) \textcolor{red}{etilciclobutano}.
    \end{solu}
\end{exercise}

\begin{exercise}
    \(\star\star\) Formula o nombra: a) but-2-eno; b) 2-metilbut-1-eno; c) ciclopenteno; e) \ch{CH3-CH=CH-CH=CH-CH3};
    \[ \text{d)}\ \chemfig{CH_2=C(-[6]CH_3)-C(-[6]CH_3)=CH-CH_3} \qquad\qquad \text{f)}\ \chemfig{*6(=(-CH_2-CH_3)-----)} \]
    \begin{solu}
        a) \textcolor{red}{\ch{CH3-CH=CH-CH3}}; b) \textcolor{red}{\chemfig{CH_2=C(-[6]CH_3)-CH_2-CH_3}}; c) \textcolor{red}{\chemfig{*5(=----)}};
        d) \textcolor{red}{2,3-dimetilpenta-1,3-dieno}; e) \textcolor{red}{hexa-2,4-dieno}; f) \textcolor{red}{1-etilciclohexeno}.
    \end{solu}
\end{exercise}

\begin{exercise}
    \(\star\star\) Formula o nombra: a) but-2-ino; b) hexa-1,3-diino; c) 4,4-dimetilpent-2-ino; d) \ch{CH+C-CH2-C+CH};
    e) \ch{CH+C-C+C-CH2-CH2-CH3};
    \[ \text{f)}\ \chemfig{CH~C-CH(-[6]CH_3)-CH_2-CH_3} \]
    \begin{solu}
        a) \textcolor{red}{\ch{CH3-C+C-CH3}}; b) \textcolor{red}{\ch{CH+C-C+C-CH2-CH3}};
        c) \textcolor{red}{\chemfig{CH_3-C~C-C(-[2]CH_3)(-[6]CH_3)-CH_3}}; d) \textcolor{red}{penta-1,4-diino};
        e) \textcolor{red}{hepta-1,3-diino}; f) \textcolor{red}{3-metilpent-1-ino}.
    \end{solu}
\end{exercise}

\begin{exercise}
    \(\star\star\) Formula o nombra: a) tolueno; b) 1,4-dimetilbenceno; c) 1-etil-3-propilbenceno; d) 1,2-dietilbenceno;
    \[ \text{e)}\ \chemfig{*6(-=-=(-CH_3)-(-Cl)=)} \qquad\qquad \text{f)}\ \chemfig{*6(-=-=(-CH_2-CH_2-CH_3)-=)} \]
    \begin{solu}
        a) \textcolor{red}{\chemfig{*6(-=-=(-CH_3)-=)}}; b) \textcolor{red}{\chemfig{*6(-=(-CH_3)-=-(-CH_3)=)}};
        c) \textcolor{red}{\chemfig{*6(-=(-CH_2-CH_3)-=(-CH_2-CH_2-CH_3)-=)}};
        d) \textcolor{red}{\chemfig{*6(-=-=(-CH_2-CH_3)-(-CH_2-CH_3)=)}};
        e) \textcolor{red}{1-cloro-2-metilbenceno} (\emph{o}-clorotolueno); f) \textcolor{red}{propilbenceno}.
    \end{solu}
\end{exercise}

\begin{exercise}
    \(\star\star\) Formula o nombra estos hidrocarburos y derivados halogenados: a) 2-metilpent-1-en-4-ino; b)
    5-metilhepta-1,3-dien-6-ino; c) 3-cloroprop-1-eno; d) 1-bromo-2-clorociclopentano; e) \ch{CH3-CHBr-CH2-CH3};
    \[ \text{f)}\ \chemfig{*6(-=(-Br)-=(-CH_3)-=)} \]
    \begin{solu}
        a) \textcolor{red}{\chemfig{CH_2=C(-[6]CH_3)-CH_2-C~CH}}; b) \textcolor{red}{\chemfig{CH_2=CH-CH=CH-CH(-[6]CH_3)-C~CH}};
        c) \textcolor{red}{\ch{CH2=CH-CH2Cl}}; d) \textcolor{red}{\chemfig{*5(-(-Br)-(-Cl)---)}}; e) \textcolor{red}{2-bromobutano};
        f) \textcolor{red}{1-bromo-3-metilbenceno}.
    \end{solu}
\end{exercise}

\begin{exercise}
    \(\star\star\) Un hornillo de acampada quema \qty{250}{\gram} de butano. a) Ajusta la reacción. b) ¿Qué volumen de oxígeno,
    medido a \qty{20}{\degreeCelsius} y \qty{1}{\atm}, se consume? c) ¿Y de aire, con un 21\,\% de \ch{O2} en volumen?
    \begin{solu}
        a) \textcolor{red}{\ch{2 C4H10 + 13 O2 -> 8 CO2 + 10 H2O}}. b) \textcolor{blue}{\(n(\ch{C4H10}) = 250/58 = \qty{4.31}{\mol}\);
        \(n(\ch{O2}) = 4{,}31 \cdot 13/2 = \qty{28.0}{\mol}\).} \(V = \dfrac{28{,}0 \cdot 0{,}082 \cdot 293}{1} =
        \textcolor{red}{\qty{673}{\litre}}\). c) \(673/0{,}21 = \textcolor{red}{\qty{3.2e3}{\litre}}\) (unos \qty{3.2}{\metre\cubed}).
    \end{solu}
\end{exercise}

\begin{exercise}
    \(\star\star\) El propino reacciona con hidrógeno en presencia de un catalizador y se transforma en propano. a) Escribe la
    reacción. ¿De qué tipo es? b) ¿Qué volumen de \ch{H2}, en condiciones normales, se necesita para hidrogenar \qty{10}{\gram} de
    propino?
    \begin{solu}
        a) \textcolor{red}{\ch{CH+C-CH3 + 2 H2 -> CH3-CH2-CH3}}: \textcolor{red}{adición}. b) \textcolor{blue}{\(n = 10/40 =
        \qty{0.25}{\mol}\) de propino; \(n(\ch{H2}) = \qty{0.50}{\mol}\).} \(V = 0{,}50 \cdot 22{,}4 = \textcolor{red}{\qty{11.2}{\litre}}\).
    \end{solu}
\end{exercise}

\subsection{Compuestos oxigenados}

\begin{exercise}
    \(\star\star\) Formula o nombra estos alcoholes y éteres: a) pentano-1,4-diol; b) 3-metilpent-3-en-2-ol; c) etil fenil éter;
    e) \ch{CH3-CH2-CH2-O-CH2-CH2-CH3};
    \[ \text{d)}\ \chemfig{CH_3-C(-[2]OH)(-[6]CH_3)-CH_2-CH_3} \qquad\qquad \text{f)}\ \chemfig{*6(-=-=(-OH)-(-CH_3)=)} \]
    \begin{solu}
        a) \textcolor{red}{\ch{CH2OH-CH2-CH2-CHOH-CH3}}; b) \textcolor{red}{\chemfig{CH_3-CH(-[6]OH)-C(-[6]CH_3)=CH-CH_3}};
        c) \textcolor{red}{\ch{C6H5-O-CH2-CH3}} (etoxibenceno); d) \textcolor{red}{2-metilbutan-2-ol};
        e) \textcolor{red}{dipropil éter, 1-propoxipropano}; f) \textcolor{red}{2-metilfenol}.
    \end{solu}
\end{exercise}

\begin{exercise}
    \(\star\star\) Formula o nombra: a) 2,2-dimetilpropanal; b) hexan-3-ona; c) ácido 3-metilbut-2-enoico; d) propanoato de etilo;
    e) pentanoato de metilo; f) \ch{CH3-CHCl-CO-CH3}; g) \ch{C6H5-COO-CH2-CH3}; h) \chemfig{CH_3-CH(-[6]CH_3)-CH_2-CH_2-COOH};
    i) \ch{OHC-CHO}; j) \ch{CH3-CH2-COONa}.
    \begin{solu}
        a) \textcolor{red}{\chemfig{CH_3-C(-[2]CH_3)(-[6]CH_3)-CHO}}; b) \textcolor{red}{\ch{CH3-CH2-CO-CH2-CH2-CH3}};
        c) \textcolor{red}{\chemfig{CH_3-C(-[6]CH_3)=CH-COOH}}; d) \textcolor{red}{\ch{CH3-CH2-COO-CH2-CH3}};
        e) \textcolor{red}{\ch{CH3-(CH2)3-COO-CH3}}; f) \textcolor{red}{3-clorobutan-2-ona}; g) \textcolor{red}{benzoato de etilo};
        h) \textcolor{red}{ácido 4-metilpentanoico}; i) \textcolor{red}{etanodial}; j) \textcolor{red}{propanoato de sodio}.
    \end{solu}
\end{exercise}

\begin{exercise}
    \(\star\star\) Ordena de menor a mayor temperatura de ebullición y justifica: propan-1-ol, butano, ácido propanoico y propanal
    (sus masas molares son parecidas: 60, 58, 74 y \qty{58}{\gram\per\mol}).
    \begin{solu}
        \textcolor{red}{Butano \(<\) propanal \(<\) propan-1-ol \(<\) ácido propanoico}. El butano es apolar (solo fuerzas de
        dispersión); el propanal es polar (dipolo-dipolo); el propanol forma enlaces de hidrógeno, y el ácido forma enlaces de
        hidrógeno más fuertes y por partida doble (dímeros). (Valores reales: \(-0{,}5\); 48; 97 y \qty{141}{\degreeCelsius}.)
    \end{solu}
\end{exercise}

\begin{exercise}
    \(\star\) Escribe la fórmula de estos compuestos y busca una aplicación de cada uno: ácido etanoico, metanal, propanona,
    etanol y etanoato de etilo.
    \begin{solu}
        \textcolor{red}{\ch{CH3-COOH}}: vinagre, conservante. \textcolor{red}{\ch{HCHO}}: formol, conservación de muestras,
        resinas. \textcolor{red}{\ch{CH3-CO-CH3}}: quitaesmaltes, disolvente. \textcolor{red}{\ch{CH3-CH2OH}}: bebidas,
        desinfectante, combustible. \textcolor{red}{\ch{CH3-COO-CH2-CH3}}: disolvente de pegamentos y lacas, aromas.
    \end{solu}
\end{exercise}

\subsection{Compuestos nitrogenados}

\begin{exercise}
    \(\star\star\) Formula o nombra: a) 3-metilpentano-2,4-diamina; b) N,N-dietilpropan-1-amina; c) pentanamida; d)
    N-etil-N-metilpropanamida; e) \ch{CH3-CH2-CO-NH-CH3}; f) \ch{CH3-CH(NH2)-CH3}; g) \ch{N(CH3)3}; h) \ch{CH3-CH2-CH2-C+N};
    \[ \text{i)}\ \chemfig{*6(-----(-NH_2)-)} \]
    \begin{solu}
        a) \textcolor{red}{\chemfig{CH_3-CH(-[6]NH_2)-CH(-[6]CH_3)-CH(-[6]NH_2)-CH_3}}; b) \textcolor{red}{\ch{CH3-CH2-CH2-N(CH2-CH3)2}};
        c) \textcolor{red}{\ch{CH3-(CH2)3-CONH2}}; d) \textcolor{red}{\ch{CH3-CH2-CO-N(CH3)-CH2-CH3}};
        e) \textcolor{red}{N-metilpropanamida}; f) \textcolor{red}{propan-2-amina}; g) \textcolor{red}{N,N-dimetilmetanamina}
        (trimetilamina); h) \textcolor{red}{butanonitrilo}; i) \textcolor{red}{ciclohexanamina}.
    \end{solu}
\end{exercise}

\subsection{Compuestos con varios grupos funcionales}

\begin{exercise}
    \(\star\star\star\) Formula o nombra: a) 5-aminohexan-2-ona; b) ácido 2-metil-3-oxobutanoico; c) 3-hidroxibutanamida;
    d) 3-aminopropan-1-ol; e) \ch{OHC-CH2-CHOH-CH3}; f) \ch{HOOC-CH2-CO-COOH}; g) \ch{CH3-CO-CH2-CH2-COOH};
    h) \ch{H2N-CH2-COOH}.
    \begin{solu}
        a) \textcolor{red}{\ch{CH3-CO-CH2-CH2-CH(NH2)-CH3}}; b) \textcolor{red}{\chemfig{CH_3-CO-CH(-[6]CH_3)-COOH}};
        c) \textcolor{red}{\ch{CH3-CHOH-CH2-CONH2}}; d) \textcolor{red}{\ch{H2N-CH2-CH2-CH2OH}}; e) \textcolor{red}{3-hidroxibutanal};
        f) \textcolor{red}{ácido 2-oxobutanodioico}; g) \textcolor{red}{ácido 4-oxopentanoico}; h) \textcolor{red}{ácido
        aminoetanoico} (glicina).
    \end{solu}
\end{exercise}

\subsection{Isomería}

\begin{exercise}
    \(\star\star\) Escribe y nombra todos los isómeros estructurales de fórmula \ch{C4H9Cl}. ¿Cuál de ellos tiene un carbono quiral?
    \begin{solu}
        \textcolor{red}{1-clorobutano, \ch{CH2Cl-CH2-CH2-CH3}; 2-clorobutano, \ch{CH3-CHCl-CH2-CH3}; 1-cloro-2-metilpropano,
        \ch{CH2Cl-CH(CH3)-CH3}; 2-cloro-2-metilpropano, \ch{(CH3)3CCl}.} Tiene carbono quiral \textcolor{red}{el 2-clorobutano}
        (C2: H, Cl, \ch{CH3} y \ch{C2H5}).
    \end{solu}
\end{exercise}

\begin{exercise}
    \(\star\star\) a) Escribe dos isómeros de función de fórmula \ch{C4H8O}. b) Escribe dos isómeros de posición de fórmula
    \ch{C5H10O} que sean cetonas.
    \begin{solu}
        a) \textcolor{red}{Butanal, \ch{CH3-CH2-CH2-CHO}, y butanona, \ch{CH3-CO-CH2-CH3}} (también 2-metilpropanal).
        b) \textcolor{red}{Pentan-2-ona, \ch{CH3-CO-CH2-CH2-CH3}, y pentan-3-ona, \ch{CH3-CH2-CO-CH2-CH3}}.
    \end{solu}
\end{exercise}

\begin{exercise}
    \(\star\star\) Indica qué tipo de isomería hay entre: a) propan-1-ol y propan-2-ol; b) ácido propanoico y etanoato de metilo;
    c) pentano y 2,2-dimetilpropano; d) \emph{cis}- y \emph{trans}-but-2-eno; e) but-1-ino y buta-1,3-dieno; f) los dos
    enantiómeros del ácido láctico.
    \begin{solu}
        a) \textcolor{red}{De posición}; b) \textcolor{red}{de función}; c) \textcolor{red}{de cadena}; d) \textcolor{red}{geométrica}
        (cis-trans); e) \textcolor{red}{de función} (un triple enlace frente a dos dobles; ambos son \ch{C4H6}); f)
        \textcolor{red}{óptica}.
    \end{solu}
\end{exercise}

\begin{exercise}
    \(\star\star\) ¿Cuáles presentan isomería geométrica y cuáles óptica? a) Pent-2-eno; b) 3-metilpent-2-eno; c) 2-bromobutano; d)
    ácido 2-aminopropanoico; e) 2-metilbutan-2-ol; f) but-1-eno.
    \begin{solu}
        \textcolor{red}{Geométrica: a) y b)} (cada carbono del doble enlace tiene dos sustituyentes distintos).
        \textcolor{red}{Óptica: c) y d)} (su C2 tiene cuatro sustituyentes distintos). e) No: el C2 tiene dos metilos. f) No: el C1
        tiene dos H.
    \end{solu}
\end{exercise}

\subsection{Importancia de los compuestos del carbono}

\begin{exercise}
    \(\star\) Escribe el monómero y un fragmento del polímero: a) PVC (policloruro de vinilo); b) teflón (politetrafluoroeteno);
    c) poliestireno, cuyo monómero es el etenilbenceno (estireno).
    \begin{solu}
        a) \textcolor{red}{\ch{CH2=CHCl}; \ch{-CH2-CHCl-CH2-CHCl-}}. b) \textcolor{red}{\ch{CF2=CF2}; \ch{-CF2-CF2-CF2-CF2-}}.
        c) \textcolor{red}{\ch{CH2=CH-C6H5}; \ch{-CH2-CH(C6H5)-CH2-CH(C6H5)-}}.
    \end{solu}
\end{exercise}

\begin{exercise}
    \(\star\star\) En la respiración celular, la glucosa se oxida: \ch{C6H12O6 + O2 -> CO2 + H2O}. a) Ajusta la ecuación.
    b) ¿Qué volumen de \ch{CO2}, medido a \qty{37}{\degreeCelsius} y \qty{1}{\atm}, se produce al oxidar \qty{45}{\gram} de glucosa?
    \begin{solu}
        a) \textcolor{red}{\ch{C6H12O6 + 6 O2 -> 6 CO2 + 6 H2O}}. b) \textcolor{blue}{\(n = 45/180 = \qty{0.25}{\mol}\);
        \(n(\ch{CO2}) = \qty{1.5}{\mol}\).} \(V = 1{,}5 \cdot 0{,}082 \cdot 310 = \textcolor{red}{\qty{38}{\litre}}\).
    \end{solu}
\end{exercise}

\begin{exercise}
    \(\star\) Indica a qué tipo de biomolécula pertenece cada sustancia y su función principal: almidón, colágeno, hemoglobina,
    triglicéridos, celulosa e insulina.
    \begin{solu}
        \textcolor{red}{Almidón: glúcido, reserva de energía en las plantas. Colágeno: proteína estructural. Hemoglobina: proteína
        de transporte de \ch{O2}. Triglicéridos: lípidos, reserva de energía y aislamiento. Celulosa: glúcido estructural (pared
        de las células vegetales). Insulina: proteína con función hormonal.}
    \end{solu}
\end{exercise}

%========================================================================================
\section{Ejercicios de ampliación}\label{sec:ud7-ampliacion}
%========================================================================================

\begin{exercise}
    \(\star\star\star\) Al quemar \qty{1.76}{\gram} de un compuesto de C, H y O se forman \qty{3.52}{\gram} de \ch{CO2} y
    \qty{1.44}{\gram} de \ch{H2O}. La densidad de su vapor a \qty{100}{\degreeCelsius} y \qty{1}{\atm} es \qty{2.87}{\gram\per\litre}.
    a) Halla su fórmula molecular. b) Escribe y nombra cuatro isómeros que sean ácidos o ésteres. c) Si tiene olor afrutado y se
    obtiene a partir de etanol y ácido acético, ¿cuál es?
    \begin{solu}
        a) \textcolor{blue}{\(n(\ch{C}) = 3{,}52/44 = \qty{0.080}{\mol}\) (\qty{0.96}{\gram}); \(n(\ch{H}) = 2 \cdot 1{,}44/18 =
        \qty{0.16}{\mol}\) (\qty{0.16}{\gram}); \(m(\ch{O}) = 1{,}76 - 1{,}12 = \qty{0.64}{\gram}\) (\qty{0.040}{\mol}): \ch{C2H4O}
        (44). \(M = d\,R\,T/p = 2{,}87 \cdot 0{,}082 \cdot 373 = \qty{88}{\gram\per\mol}\).} \textcolor{red}{\ch{C4H8O2}}.
        b) \textcolor{red}{Ácido butanoico, ácido 2-metilpropanoico, etanoato de etilo, propanoato de metilo} (y metanoato de propilo
        o de isopropilo). c) \textcolor{red}{Etanoato de etilo}, \ch{CH3-COO-CH2-CH3}.
    \end{solu}
\end{exercise}

\begin{exercise}
    \(\star\star\star\) \qty{10}{\litre} de una mezcla de metano y propano se queman por completo y producen \qty{22}{\litre} de
    \ch{CO2}, medidos en las mismas condiciones. a) Halla la composición de la mezcla en volumen. b) ¿Qué volumen de \ch{O2} se ha
    consumido?
    \begin{solu}
        \textcolor{blue}{\ch{CH4 + 2 O2 -> CO2 + 2 H2O}; \ch{C3H8 + 5 O2 -> 3 CO2 + 4 H2O}. Con \(x\) litros de metano:
        \(x + 3(10 - x) = 22\), \(x = 4\).} a) \textcolor{red}{\qty{4}{\litre} de metano (40\,\%) y \qty{6}{\litre} de propano (60\,\%)}.
        b) \(2 \cdot 4 + 5 \cdot 6 = \textcolor{red}{\qty{38}{\litre}}\) de \ch{O2}.
    \end{solu}
\end{exercise}

\begin{exercise}
    \(\star\star\star\) Una muestra de \qty{2.80}{\gram} de un alqueno con un solo doble enlace reacciona exactamente con
    \qty{8.00}{\gram} de bromo, \ch{Br2}, que se adiciona al doble enlace. a) Halla la fórmula molecular del alqueno. b) Escribe
    sus isómeros de cadena abierta. c) Escribe y nombra el producto si el alqueno es el but-2-eno.
    \begin{solu}
        a) \textcolor{blue}{\(n(\ch{Br2}) = 8{,}00/160 = \qty{0.0500}{\mol} = n(\text{alqueno})\); \(M = 2{,}80/0{,}0500 =
        \qty{56}{\gram\per\mol}\); \(14n = 56\), \(n = 4\).} \textcolor{red}{\ch{C4H8}}. b) \textcolor{red}{But-1-eno,
        \emph{cis}- y \emph{trans}-but-2-eno y metilpropeno}. c) \textcolor{red}{\ch{CH3-CHBr-CHBr-CH3}, 2,3-dibromobutano}.
    \end{solu}
\end{exercise}

\begin{exercise}
    \(\star\star\star\) Nombra:
    \begin{gather*}
        \text{a)}\ \chemfig{CH_2=CH-CH(-[6]CH=[0]CH_2)-CH_2-CH_3} \qquad \text{b)}\ \chemfig{CH~C-CH_2-CH(-[6]CH_3)-CH_2OH} \\
        \text{c)}\ \chemfig{CH_3-CH_2-CH(-[6]CHO)-CH_2-CH_2-CH_3}
    \end{gather*}
    \begin{solu}
        a) \textcolor{blue}{Todas las cadenas tienen 5 C; se elige la que contiene los dos dobles enlaces.}
        \textcolor{red}{3-etilpenta-1,4-dieno}. b) \textcolor{blue}{La cadena contiene el \ch{-OH} (C1).}
        \textcolor{red}{2-metilpent-4-in-1-ol}. c) \textcolor{blue}{El carbono del \ch{-CHO} forma parte de la cadena principal (C1);
        la más larga que lo contiene tiene 5 C.} \textcolor{red}{2-etilpentanal}.
    \end{solu}
\end{exercise}

\begin{exercise}
    \(\star\star\star\) Formula: a) ácido 4-amino-3-hidroxibutanoico; b) ácido 2-amino-3-hidroxipropanoico (serina); c)
    4-hidroxi-4-metilpentan-2-ona; d) 2-aminobenzoato de metilo (aroma de uva).
    \begin{solu}
        a) \textcolor{red}{\ch{H2N-CH2-CHOH-CH2-COOH}}; b) \textcolor{red}{\ch{HOCH2-CH(NH2)-COOH}};
        c) \textcolor{red}{\chemfig{CH_3-CO-CH_2-C(-[2]OH)(-[6]CH_3)-CH_3}}; d) \textcolor{red}{\chemfig{*6(-=-=(-COO-CH_3)-(-NH_2)=)}}.
    \end{solu}
\end{exercise}

\begin{exercise}
    \(\star\star\star\) La aspirina (ácido acetilsalicílico, \ch{C9H8O4}) se obtiene a partir del ácido salicílico (\ch{C7H6O3}) en
    proporción 1 : 1. Un laboratorio obtiene \qty{2.88}{\gram} de aspirina a partir de \qty{2.76}{\gram} de ácido salicílico. ¿Cuál
    es el rendimiento?
    \begin{solu}
        \textcolor{blue}{\(M(\ch{C7H6O3}) = 138\), \(M(\ch{C9H8O4}) = \qty{180}{\gram\per\mol}\). \(n = 2{,}76/138 = \qty{0.0200}{\mol}\);
        masa teórica \(0{,}0200 \cdot 180 = \qty{3.60}{\gram}\).} \(\dfrac{2{,}88}{3{,}60} \cdot 100 = \textcolor{red}{\qty{80}{\percent}}\).
    \end{solu}
\end{exercise}

\begin{exercise}
    \(\star\star\) En la síntesis de Wöhler, el cianato de amonio, \ch{NH4OCN}, se transforma en urea, \ch{(NH2)2CO}. a) ¿Qué masa
    de urea se obtiene, como máximo, a partir de \qty{15}{\gram} de cianato de amonio? b) ¿Qué relación hay entre las dos
    sustancias? c) ¿Qué grupo funcional orgánico tiene la urea?
    \begin{solu}
        a) \textcolor{red}{\qty{15}{\gram}}: las dos sustancias tienen la misma fórmula, \ch{CH4N2O}, y la misma masa molar (60).
        b) \textcolor{red}{Son isómeros}. c) \textcolor{red}{Amida} (es la diamida del ácido carbónico).
    \end{solu}
\end{exercise}

\begin{exercise}
    \(\star\star\) La fermentación alcohólica de la glucosa es: \ch{C6H12O6 -> 2 CH3-CH2OH + 2 CO2}. Con \qty{1.0}{\kilo\gram} de
    glucosa, calcula: a) la masa y el volumen de etanol que se obtienen (\(d = \qty{0.79}{\gram\per\mL}\)); b) el volumen de
    \ch{CO2} desprendido a \qty{25}{\degreeCelsius} y \qty{1}{\atm}.
    \begin{solu}
        \textcolor{blue}{\(n(\text{glucosa}) = 1000/180 = \qty{5.56}{\mol}\); \(n(\text{etanol}) = n(\ch{CO2}) = \qty{11.1}{\mol}\).}
        a) \(m = 11{,}1 \cdot 46 = \textcolor{red}{\qty{511}{\gram}}\); \(V = 511/0{,}79 = \textcolor{red}{\qty{647}{\mL}}\).
        b) \(V = 11{,}1 \cdot 0{,}082 \cdot 298 = \textcolor{red}{\qty{271}{\litre}}\).
    \end{solu}
\end{exercise}

\begin{exercise}
    \(\star\star\star\) a) ¿Cuántos carbonos quirales tiene el 2,3-dibromobutano? b) ¿Y la forma de cadena abierta de la glucosa,
    \ch{CHO-CHOH-CHOH-CHOH-CHOH-CH2OH}?
    \begin{solu}
        a) \textcolor{red}{Dos}, el C2 y el C3 (cada uno unido a H, Br, \ch{CH3} y al otro carbono). b) \textcolor{red}{Cuatro}, del C2
        al C5; el C1 (\ch{-CHO}) y el C6 (\ch{-CH2OH}) no lo son. Por eso existen \(2^4 = 16\) estereoisómeros de esa fórmula.
    \end{solu}
\end{exercise}

%========================================================================================
\section{Autoevaluación}\label{sec:ud7-autoevaluacion}
%========================================================================================

\begin{exercise}
    ¿Por qué el carbono forma tantos compuestos?
    \begin{solu}
        \textcolor{red}{Porque es tetravalente (forma cuatro enlaces covalentes), sus enlaces \ch{C-C} son muy estables y le
        permiten formar cadenas largas, ramificadas o cerradas, con enlaces sencillos, dobles o triples}, y además se une a H, O, N y
        halógenos.
    \end{solu}
\end{exercise}

\begin{exercise}
    Nombra y clasifica en alcanos, alquenos, alquinos o aromáticos: a) \ch{CH3-(CH2)4-CH3}; b) \ch{CH3-CH=CH-CH3};
    c) \ch{CH+C-CH2-CH3}; d) \ch{C6H5-CH2-CH3}.
    \begin{solu}
        a) \textcolor{red}{Hexano, alcano}; b) \textcolor{red}{but-2-eno, alqueno}; c) \textcolor{red}{but-1-ino, alquino};
        d) \textcolor{red}{etilbenceno, aromático}.
    \end{solu}
\end{exercise}

\begin{exercise}
    Formula: a) octano; b) hexa-1,3-dieno; c) but-1-ino; d) etilciclohexano; e) 2,3-dicloropentano; f) etanol; g) propanal;
    h) metil propil éter; i) 2,2-dimetilbutanal; j) 4-metilpent-3-en-2-ona; k) ácido pentanoico; l) etanoato de metilo;
    m) N-metilpropanamida; n) N,N-dimetiletanamina; ñ) anilina.
    \begin{solu}
        a) \textcolor{red}{\ch{CH3-(CH2)6-CH3}}; b) \textcolor{red}{\ch{CH2=CH-CH=CH-CH2-CH3}}; c) \textcolor{red}{\ch{CH+C-CH2-CH3}};
        d) \textcolor{red}{\chemfig{*6(-----(-CH_2-CH_3)-)}}; e) \textcolor{red}{\ch{CH3-CHCl-CHCl-CH2-CH3}};
        f) \textcolor{red}{\ch{CH3-CH2OH}}; g) \textcolor{red}{\ch{CH3-CH2-CHO}}; h) \textcolor{red}{\ch{CH3-O-CH2-CH2-CH3}};
        i) \textcolor{red}{\chemfig{CH_3-CH_2-C(-[2]CH_3)(-[6]CH_3)-CHO}}; j) \textcolor{red}{\chemfig{CH_3-CO-CH=C(-[6]CH_3)-CH_3}};
        k) \textcolor{red}{\ch{CH3-(CH2)3-COOH}}; l) \textcolor{red}{\ch{CH3-COO-CH3}}; m) \textcolor{red}{\ch{CH3-CH2-CO-NH-CH3}};
        n) \textcolor{red}{\ch{CH3-CH2-N(CH3)2}}; ñ) \textcolor{red}{\ch{C6H5-NH2}}.
    \end{solu}
\end{exercise}

\begin{exercise}
    Nombra: a) \ch{CH3-(CH2)5-CH3}; b) \ch{CH3-CHOH-CHOH-CH3}; c) \ch{CH+C-C+C-CH3}; d) \ch{CH3-CH2-CHBr-CHBr-CH3};
    e) \ch{CH3-CH2-CO-CH2-CH3}; f) \ch{CH3-CH2-COO-CH2-CH3}; g) \ch{CH3-CH2-NH-CH3}; h) \ch{HCOOH}; i) \ch{CH3-CONH2}.
    \begin{solu}
        a) \textcolor{red}{Heptano}; b) \textcolor{red}{butano-2,3-diol}; c) \textcolor{red}{penta-1,3-diino};
        d) \textcolor{red}{2,3-dibromopentano}; e) \textcolor{red}{pentan-3-ona}; f) \textcolor{red}{propanoato de etilo};
        g) \textcolor{red}{N-metiletanamina}; h) \textcolor{red}{ácido metanoico}; i) \textcolor{red}{etanamida}.
    \end{solu}
\end{exercise}

\begin{exercise}
    Escribe y nombra dos isómeros de cadena de fórmula \ch{C5H12}.
    \begin{solu}
        Por ejemplo, \textcolor{red}{pentano, \ch{CH3-(CH2)3-CH3}, y 2-metilbutano, \chemfig{CH_3-CH(-[6]CH_3)-CH_2-CH_3}}.
    \end{solu}
\end{exercise}

\begin{exercise}
    Escribe y nombra dos isómeros de posición de fórmula \ch{C4H10O}.
    \begin{solu}
        \textcolor{red}{Butan-1-ol, \ch{CH3-CH2-CH2-CH2OH}, y butan-2-ol, \ch{CH3-CHOH-CH2-CH3}}.
    \end{solu}
\end{exercise}

\begin{exercise}
    Escribe y nombra dos isómeros de función de fórmula \ch{C3H8O}.
    \begin{solu}
        \textcolor{red}{Propan-1-ol, \ch{CH3-CH2-CH2OH}, y metoxietano (etil metil éter), \ch{CH3-O-CH2-CH3}}.
    \end{solu}
\end{exercise}

\fussy
